\scanpage{521}
\mspage{432}
% B01
et même dans $E_{s\pi}$, puisque $E_{s\pi}$ est stable dans $E_s$ par facteurs directs. Cela prouve que $x\in\mathrm{Ob}\,E(\pi)$, qed.

% B02
On peut présenter 4.15.1 sous une forme plus \unsure{sympa} :

% B03
\uline{Corollaire} 4.15.\unsure{4|9|3} : Soit $f:E\longrightarrow C$ un foncteur de $\unsure{\mathcal U|\mathscr U|U}$-catégories. \inserted{Les} conditions suivantes sont équivalentes :

% B04
a) $\forall$ foncteur $C'\longrightarrow C$, avec $C'$ petite et sans projecteurs, $E'\overset{\mathrm{déf}}{=}E\times_C C'$ est accessible, et $f':\underline E'\unsure{\mathord{-}|\longrightarrow}C'$ est accessible.

% B05
b) Comme a), mais avec $C'=\unsure{\mathbb{\Delta}|\Delta}^{1}$, en d'autres termes c'est la condition $H$) p. 409 : pour tte $\lambda\in\mathrm{Fl}\,C$, $E_\lambda\unsure{=|\Rightarrow}E\times_C(\mathbb{\Delta}1,i_\lambda)$ est accessible, et $E_\lambda\longrightarrow\mathbb{\Delta}^1$ est accessible.

% B06
c) La condition $H'=H1+H2+H3$, explicitée dans l'énoncé de 4.15.1.

% B07
\uline{Démonstration} Elle se fait circulaire
% B08
\[a\Longrightarrow b\Longrightarrow c\Longrightarrow a,\]
% B09
l'implication $a\Longrightarrow b$ étant tautologique, $b\Longrightarrow c$ l'étant presque (cf \deleted{\unsure{démonstration}} de 4.14.24 $a\Longrightarrow b\Longrightarrow c$). \deleted{\unsure{Il}}En fait les conditions $b$ et $c$ sont des conditions sur les catégories $E_\lambda\xrightarrow{f_\lambda}\mathbb{\Delta}^1$ pour $\lambda\in\mathrm{Fl}\,C$, \unsure{aussi|ainsi|aus|donc} stables par chgt de base. Pour prouver

% archive sheet number: 520
